% Standalone problem document, generated from the adjacent README.md.
% Compile with XeLaTeX. Mathematical content is maintained in Markdown.
\documentclass[11pt,a4paper]{article}
\usepackage[fontsize=10.5pt]{scrextend}
\usepackage[margin=22mm,headheight=15pt,headsep=9mm,footskip=11mm]{geometry}
\usepackage{fontspec}
\setmainfont{texgyrepagella}[Extension=.otf,UprightFont=*-regular,BoldFont=*-bold,ItalicFont=*-italic,BoldItalicFont=*-bolditalic]
\setsansfont{texgyreheros}[Extension=.otf,UprightFont=*-regular,BoldFont=*-bold,ItalicFont=*-italic,BoldItalicFont=*-bolditalic]
\setmonofont{texgyrecursor}[Extension=.otf,UprightFont=*-regular,BoldFont=*-bold,ItalicFont=*-italic,BoldItalicFont=*-bolditalic,Scale=MatchLowercase]
\usepackage{amsmath,amssymb}
\usepackage{unicode-math}
\setmathfont{texgyrepagella-math.otf}
\usepackage{xcolor,microtype,enumitem,needspace,fancyhdr}
\usepackage{bookmark}
\definecolor{ink}{HTML}{17344B}
\definecolor{accent}{HTML}{197D80}
\definecolor{muted}{HTML}{596773}
\hypersetup{colorlinks=true,linkcolor=accent,urlcolor=accent,pdftitle={MI-24 - Schatten
norm complement involving Lin's positive definite
quantity},pdfauthor={Open Problems in Numerical Linear Algebra}}
\urlstyle{same}
\setlength{\parindent}{0pt}
\setlength{\parskip}{5pt plus 1pt minus 1pt}
\linespread{1.04}
\setlength{\emergencystretch}{3em}
\widowpenalty=10000
\clubpenalty=10000
\displaywidowpenalty=10000
\allowdisplaybreaks[1]
\setlist{nosep,leftmargin=1.5em}
\providecommand{\tightlist}{\setlength{\itemsep}{0pt}\setlength{\parskip}{0pt}}
\providecommand{\pandocbounded}[1]{#1}
\pagestyle{fancy}
\fancyhf{}
\fancyhead[L]{\sffamily\footnotesize\color{muted}OPEN PROBLEMS IN NUMERICAL LINEAR ALGEBRA}
\fancyhead[R]{\sffamily\bfseries\footnotesize\color{accent}MI-24}
\fancyfoot[L]{\parbox[b]{0.9\textwidth}{\sffamily\fontsize{8}{10}\selectfont\color{muted}Editorial ratings; a bounded search cannot certify that a problem remains unsolved.\\Literature check: 2026-09-10}}
\fancyfoot[R]{\sffamily\footnotesize\color{muted}\thepage}
\renewcommand{\headrulewidth}{0.3pt}
\renewcommand{\headrule}{\hbox to\headwidth{\color{accent}\leaders\hrule height \headrulewidth\hfill}}
\makeatletter
\renewcommand{\section}{\Needspace{5\baselineskip}\@startsection{section}{1}{0pt}{12pt plus 2pt minus 2pt}{4pt}{\sffamily\bfseries\large\color{ink}}}
\renewcommand{\subsection}{\Needspace{4\baselineskip}\@startsection{subsection}{2}{0pt}{9pt plus 1pt minus 1pt}{3pt}{\sffamily\bfseries\normalsize\color{ink}}}
\makeatother
\setcounter{secnumdepth}{0}
\begin{document}
{\sffamily\small\color{muted}Matrix inequalities and norms\par}
\vspace{3pt}
{\raggedright\hyphenpenalty=10000\exhyphenpenalty=10000\sffamily\bfseries\fontsize{21}{24}\selectfont\color{ink}Schatten
norm complement involving Lin's positive definite quantity\par}
\vspace{7pt}
{\sffamily\small\textbf{Difficulty:} challenging\quad\textbf{Importance:} interesting
to specialist\par}
{\sffamily\small\bfseries\color{accent}Status: Partially resolved\par}
\vspace{4pt}
{\color{accent}\hrule height 0.8pt}
\vspace{6pt}
\textbf{Rating rationale:} Extending the proved trace and Frobenius
cases to all Schatten norms is challenging; the comparison of these
particular means has specialist importance.

\subsection{Problem statement}\label{problem-statement}

For positive definite \(A,B\in\mathbb C^{n\times n}\), set \[
G=A^{1/2}(A^{-1/2}BA^{-1/2})^{1/2}A^{1/2},\qquad
L=A^{1/2}(B^{1/2}A^{-1}B^{1/2})^{1/2}A^{1/2}.
\] Here \(G=A\#B\), while \(L\) is the explicitly defined quantity used
by Lin; no symmetry of \(L(A,B)\) is assumed. Determine whether, for all
\(n\ge1\), all positive definite \(A,B\), and all \(1\le p\le\infty\),
\[
\|A+B+G+L\|_p\le\|A+B+2L\|_p.
\] For \(p<\infty\), \(\|X\|_p=(\operatorname{tr}|X|^p)^{1/p}\), where
\(|X|=(X^*X)^{1/2}\); for \(p=\infty\), use the largest singular value.

Matrix means are central to interpolation of positive definite data. The
conjecture supplies a norm comparison between two concrete sums of
matrix functions; eigenvalue comparison between \(G\) and \(L\) alone is
not the statement being requested.

\subsection{References}\label{references}

\begin{enumerate}
\def\labelenumi{\arabic{enumi}.}
\tightlist
\item
  M. M. Ghabries, H. Abbas, B. Mourad and A. Assi, \emph{New
  log-majorization results concerning eigenvalues and singular values
  and a complement of a norm inequality}, preprint (2021), definition in
  §4, p.~11; Conjecture 4.1, p.~13.
  \href{https://arxiv.org/pdf/2105.13356}{PDF};
  \href{https://doi.org/10.1080/03081087.2022.2059050}{journal article}.
\item
  M. M. Ghabries, \emph{Contributions to Matrix Inequalities and Some
  Applications}, PhD thesis (2022), final Open Problems, Problem 5,
  pp.~109--110.
  \href{https://www.researchgate.net/publication/361793582_Contributions_to_Matrix_Inequalities_and_Some_Applications}{Author-uploaded
  thesis}.
\end{enumerate}

Status check (2026-09-10): the first source proves \(p=1,2\) and then
conjectures the full interval. The thesis explicitly retains that gap.
Current arXiv version v1, published metadata, the author's later July
2026 matrix-mean paper, and searches for ``Ghabries Conjecture 4.1
norm'', ``Lin complement geometric mean Schatten'', including 2025/2026,
did not reveal a proof or counterexample for the full interval. No
assertion of exhaustive coverage is made.
\end{document}
